\documentclass[12pt,a4paper]{report}

% ============================ 预加载的 LaTeX 包 ============================
\usepackage[a4paper,left=3.5cm,right=3cm,top=3cm,bottom=3cm]{geometry} % 页面边距
\usepackage[protrusion=true,expansion=false]{microtype} % XeTeX不支持字体扩展
\usepackage{setspace}                         % 设置 1.5 倍行距
\usepackage{fancyhdr}                         % 自定义页眉页脚
% amsmath/amsthm 须在 newtxmath 之前；勿加载 amssymb（与 newtxmath 的 AMS 符号重复）
\usepackage{amsmath,amsthm}                   % 数学与定理环境
\usepackage{newtxtext,newtxmath}              % Times New Roman 字体
\usepackage{hyperref}                         % 目录超链接
\hypersetup{hidelinks,colorlinks=false,linktoc=all}
\usepackage{booktabs}                         % 美化表格
\usepackage{array}                            % 处理表格对齐
\usepackage{tocbibind}  % 让目录中包含“参考文献”和“附录”
\usepackage{caption}                          % 处理表格和图片标题格式
\usepackage{xcolor}                           % 颜色支持
\usepackage{titlesec}                         % 控制标题字体大小
\usepackage{tocloft}                          % 目录格式控制
\usepackage{cite}                             % bib引用
% \usepackage{showframe}
\usepackage{etoolbox}

% ================================= 设置公式 ================================
\AtBeginEnvironment{equation}{\setstretch{1}}             % 设置公式1倍行距
\AtEndEnvironment{equation}{}                             % 公式结束恢复原样式
\numberwithin{equation}{chapter}                          % 让公式编号跟随章节
\renewcommand{\theequation}{\arabic{chapter}.\arabic{equation}}% 仅显示1.1

% ================================= 设置页脚 ================================
\pagestyle{fancy}                                         % 启用 fancyhdr
\fancyhf{}                                                % 清除默认的页眉和页脚
\fancyfoot[R]{\fontsize{12pt}{14pt}\selectfont \thepage}  % 右侧页脚页码 12pt
\renewcommand{\headrulewidth}{0pt}                        % 取消页眉下划线
\renewcommand{\footrulewidth}{0pt}                        % 取消页脚下划线
\fancypagestyle{plain}{                                   % 适配plain风格页面
    \fancyhf{} 
    \fancyfoot[R]{\fontsize{12pt}{14pt}\selectfont \thepage}
    \renewcommand{\headrulewidth}{0pt}
    \renewcommand{\footrulewidth}{0pt}
}

% ================================= 设置目录 ================================
% 修改目录标题格式（居中 + 18pt + 加粗 + 下划线）
\renewcommand{\contentsname}{}
\newcommand{\customtoc}{
    \clearpage
    \begin{center}
        \underline{\textbf{\fontsize{18pt}{18pt}\selectfont Table of Contents}}
    \end{center}
    \vspace{-1.2\baselineskip}
}
% 自定义 "Pages" 列标题，使其右对齐
\renewcommand{\cftaftertoctitle}{%
    \vspace{-1.5\baselineskip} % 减少 Table of Contents 后的额外空行
    \hfill\underline{Pages} % 右对齐
    \par\vskip-2\baselineskip % 调整间距，防止错位
}
%目录格式调整
% 一级标题 (chapter) 设置
\renewcommand{\cftchapfont}{\bfseries}                  % 章节标题加粗
\renewcommand\thechapter{\arabic{chapter}}
\renewcommand{\cftchappresnum}{Chapter\enspace}
\renewcommand\thesection{\arabic{chapter}.\arabic{section}} % 解决上行代码的乱码问题
\renewcommand{\cftchapaftersnumb}{}                     % 章节编号后不添加任何内容
\setlength{\cftchapnumwidth}{5.3em}                       % 设置编号宽度
% 二级标题（section）设置
\renewcommand{\cftsecfont}{}                            % 节标题字体
\renewcommand{\cftsecpresnum}{}                         % 节编号前缀
\renewcommand{\cftsecaftersnumb}{}                      % 节编号后缀
\setlength{\cftsecindent}{0em}                          % 节的缩进
\setlength{\cftsecnumwidth}{2em}                        % 节编号宽度
\renewcommand{\cftsecdotsep}{\cftnodots}                % 移除点线
\renewcommand{\cftsecpagefont}{\color{white}}        % 移除页码
% 间距设置
\setlength{\cftbeforechapskip}{0.5em}                   % 章节前的间距
\setlength{\cftbeforesecskip}{0.5em}                    % 节前的间距

% ========================== 统一设置 \chapter ==========================
% 设置有编号的 \chapter 显示 "Chapter 1" + 换行后的标题左对齐
\titleformat{\chapter}[display]            % "Chapter 1" 独立一行，标题换行后左对齐
  {\normalfont\Huge\bfseries}  % 不倾斜，字号 Huge，加粗
  {\chaptername\enspace\thechapter}         % 保留可见章标题，引用仅使用数字
  {0em}                                    % 章节编号和标题之间的间距
  {\raggedright}                           % 标题居左
% 设置无编号的 \chapter*{}（如摘要、附录等）格式，不带 "Chapter"
\titleformat{name=\chapter,numberless}[display]
  {\normalfont\Huge\bfseries}  % 不倾斜，字号 Huge，加粗
  {}                                       % 无编号时不显示 "Chapter 1"
  {-2.5em}                                 % 章节编号和标题之间的间距
  {\centering}                             % 标题居中
% 调整章节标题的间距
\titlespacing{\chapter}{0pt}{-2em}{0em}  % 影响有编号的章节
\titlespacing{name=\chapter,numberless}{0pt}{0em}{\baselineskip}  % 影响无编号的章节

% ========================== 统一设置 \section ==========================
% 设置有编号的 \section{}
\titleformat{\section}                        % 统一设置 section
  {\normalfont\fontsize{16pt}{16pt}\bfseries} % 12pt 字号行距，加粗
  {\thesection\enspace}                       % 章节编号（这里为空）
  {0pt}                                       % 章节前间距
  {\raggedright}                              % 章节后格式（这里保持默认）
% 设置无编号的 \section*{}
\titleformat{name=\section,numberless}[display]
  {\clearpage\normalfont\fontsize{16pt}{16pt}\bfseries}
  {}
  {-1em}
  {\centering}
\titlespacing{\section}{0pt}{0.5\baselineskip}{0.5\baselineskip}
\titlespacing{name=\section,numberless}{0pt}{0em}{1.5\baselineskip}

\global\sloppy  % 允许换行，减少溢出

% ======================= 统一设置 \subsection ========================
% 设置有编号的 \subsection{}
\titleformat{\subsection}    
  {\normalfont\fontsize{14pt}{14pt}\bfseries}  % 14pt 字号行距，加粗
  {\thesubsection\enspace}                     % 章节编号
  {0pt}                                        % 章节编号与标题间距
  {\raggedright}                               % 标题左对齐

% 调整 subsection 的间距
\titlespacing{\subsection}{0pt}{0.5\baselineskip}{0.5\baselineskip}


\begin{document}
\onehalfspacing
\setlength{\parindent}{0pt}
\setlength{\parskip}{0.3em}
\setlength{\emergencystretch}{2em}
\renewcommand{\cftsecpagefont}{\normalfont}
\renewcommand{\contentsname}{Contents}
\pagenumbering{gobble}
\begin{titlepage}
\centering
\vspace*{2cm}
{\Large\bfseries Graph-Augmented Vision Transformers\\[0.4em]
for Remote Sensing Image Scene Understanding\par}
\vspace{2cm}
{\large Dissertation working draft: Chapters 1--4\par}
\vspace{0.5cm}
{\large English chapter draft\par}
\vfill
29 October 2026\par
\vspace{0.7cm}
{\small This partial draft contains the introduction, literature review,
method and experimental setup. Formal front matter, comparative results,
discussion and conclusions remain to be integrated.}
\end{titlepage}
\pagenumbering{roman}
\tableofcontents
\clearpage
\pagenumbering{arabic}
\setcounter{chapter}{0}
\chapter{Introduction}
\label{chap:introduction}

\section{Background}

Remote sensing scene classification assigns semantic categories to aerial or
satellite image patches. It provides a means of organising image archives and
describing land use and land cover at the scene level. These descriptions can
support applications such as urban analysis and environmental monitoring,
although an image-level prediction does not locate every object or delineate
every land-cover boundary. Surveys by Zhu et al.~\cite{zhu2017deeprs} and
Cheng et al.~\cite{cheng2020scene} place scene classification within the wider
development of learned representations for remote sensing.

The target description depends on the task. A single-label benchmark assigns
one scene category to each image, even when several visual elements coexist.
NWPU-RESISC45 and the Aerial Image Dataset (AID) provide examples of this
setting~\cite{cheng2017nwpu,xia2017aid}. A multi-label task instead allows an
image to carry several land-cover labels, as in
BigEarthNet~\cite{sumbul2019bigearthnet,sumbul2021bigearthnetmm}. This
distinction changes the prediction objective and evaluation: choosing one
class differs from ranking and thresholding multiple label probabilities.
It also motivates examining regional representations in both settings,
without assuming that a useful representation for one will transfer
unchanged to the other.

Scene appearance presents a further challenge. Images within one category
may differ in scale, orientation, texture and composition, while different
categories can share visually similar components~\cite{cheng2017nwpu}.
Consequently, recognising local elements and combining their contextual
information are closely connected problems. An image representation must
retain useful evidence while reducing the many spatial features to the
compact vector required by a classifier.

Transfer learning offers a practical starting point. ImageNet established a
large natural-image resource~\cite{deng2009imagenet}, and Nogueira
et al.~\cite{nogueira2017exploiting} examined how pretrained convolutional
networks can be used for remote sensing scene classification. Domain-specific
pretraining provides another route: Seasonal Contrast uses temporal imagery
for self-supervised representation learning~\cite{manas2021seco}, while
SatMAE adapts masked pretraining to temporal and multispectral satellite
data~\cite{cong2022satmae}. These approaches illustrate that pretraining and
downstream architecture are separate design choices. This thesis holds the
recorded pretrained backbone source common within its internal comparison
and concentrates on the structure added after feature extraction.

\section{Motivation}

The shared backbone is Swin-Tiny (Swin-T), a hierarchical vision Transformer whose
shifted-window design enables interaction across neighbouring
windows~\cite{liu2021swin}. Its final tokens already contain contextual
information. The motivation for adding a graph is therefore to investigate
an additional regional inductive bias, rather than to remedy an assumed
absence of spatial relationships in the backbone. Whether that bias is
useful is an empirical question: an extra module can introduce helpful
structure, redundant computation or additional sensitivity to training.

Regional aggregation offers an intermediate representation between dense
feature tokens and a single global vector. Grouping tokens into a smaller
set of nodes makes their interactions explicit. Graph attention can then
aggregate neighbouring node features using learned
weights~\cite{velickovic2018gat}. However, this description leaves important
choices unresolved. Region membership, pooling weights, edge selection and
the return of graph features to the classifier each determine what
information is preserved and exchanged. An apparent improvement from a
complete architecture cannot automatically identify which choice caused it.

This thesis examines two grouping strategies within a Graph-Augmented Vision
Transformer (GAViT). The fixed strategy uses a deterministic spatial partition and
learns the relative importance of tokens within each region. The soft
strategy constructs image-dependent prototypes and permits contributions to
nearby anchors. Both retain a common anchor layout and explicit edge budget.
The comparison asks whether adaptive membership improves on a simple spatial
prior when the backbone and subsequent graph machinery are shared. It does
not equate attention-weighted pooling with learned segmentation, or assume
that soft assignment necessarily yields more meaningful regions.

The value of this structure also depends on its cost and consistency.
Parameter count, latency and memory describe different aspects of the added
computation. Repeated training can reveal sensitivity that a single favourable
score conceals. Visual inspection can show actual assignments, connections
and prediction changes, but cannot alone establish semantic correctness or
causal benefit. These considerations motivate evaluating the complete
accuracy--cost trade-off alongside diagnostic evidence.

\section{Research questions}

Three research questions organise the investigation:
\begin{enumerate}
 \item \textbf{RQ1: Regional graph reasoning.} Does fixed-partition GAViT
 improve classification relative to the shared Swin-T baseline, and how
 consistent is the observed difference across eligible training repetitions?
 \item \textbf{RQ2: Adaptive membership.} Does local prototype soft assignment
 improve on fixed grouping under the same task protocol, and what changes
 can be observed in regional aggregation, graph connectivity and predictions?
 \item \textbf{RQ3: Practical value and interpretation.} How do classification
 differences relate to additional computational cost, and what do successes,
 failures and visual diagnostics support about the role of the regional
 representation?
\end{enumerate}
These questions permit conditional and negative findings. A more flexible
grouping mechanism is not treated as preferable merely because it contains
additional learnable parameters.

\section{Contributions and scope}

The work makes three contributions. First, it implements a
Graph-Augmented Vision Transformer that combines attentive regional
aggregation, a sparse spatial-and-feature graph, graph attention and
region-to-token feedback. The architectural description distinguishes fixed
membership from learned pooling and distinguishes feature-based edge
selection from graph-attention weights. This makes the computational roles
of the components explicit.

Second, it formulates an image-dependent local prototype grouping mechanism
within the same architecture. Distinct normalisations define token membership
and regional pooling, while feedback uses the membership weights. This allows
the comparison to address adaptive aggregation without replacing the shared
backbone. The soft-versus-fixed
comparison concerns the resulting grouping pathway as a whole; it does not
separately identify every effect of the added projection and feedback map.

Third, it establishes an evaluation framework combining single-label and
multi-label classification, explicit model-selection rules, training
repetitions, computational measurements and qualitative diagnostics.
Calibration and final fitting are distinguished, and repeated training on
a frozen partition is distinguished from repeated random data splitting.
The contribution is the implemented design and its evidence-based
examination; consistent superiority over the baseline is a proposition to
be tested rather than an assumed outcome.

The empirical scope comprises red, green and blue (RGB) scene classification on NWPU-RESISC45,
AID and BigEarthNet-19. BigEarthNet is restricted to three optical bands in
this study, so the work does not evaluate full multispectral, radar or
multitemporal fusion. It also does not perform object detection or semantic
segmentation. Regional visualisations describe internal feature aggregation,
not independently validated object masks. Independent reconstruction of
HGTNet and DFAGCN forms a complementary comparison track; formal external
comparisons require the remaining implementation and protocol ambiguities
to be resolved.

\section{Organisation of the thesis}

The thesis is organised into seven chapters. Following this introduction,
Chapter~\ref{chap:related-work} reviews remote sensing Transformers, image graph networks and
regional grouping. Chapter~\ref{chap:method} presents GAViT and its fixed and
soft grouping variants. Chapter~\ref{chap:setup} specifies datasets,
preprocessing, model selection and evaluation. Chapter~5 reports
classification comparisons and component studies. Chapter~6 analyses
computational cost, qualitative evidence, failures and limitations.
Chapter~7 concludes the thesis and identifies directions for future work.

\input{chapters/02-literature-review}
\input{chapters/03-method}
\chapter{Experimental Setup}
\label{chap:setup}

\section{Evaluation questions}

The experiments address three questions. First, does the fixed region graph
improve classification relative to the shared Swin-T baseline? Second, does
the local prototype soft assignment improve on the fixed grouping under the
same task protocol? Third, are any observed benefits consistent across
training repetitions and sufficient to justify the additional computational
cost? Classification results, component comparisons and diagnostic
visualisations provide different kinds of evidence for these questions.

This chapter specifies the implemented internal comparison. The external
comparisons with HGTNet and DFAGCN are retained in the thesis scope, but their
common training recipe and experimental results remain separate pending
implementation materials and protocol verification. Published scores are
not treated as measurements made under the internal protocol.

\section{Datasets and fixed partitions}

\subsection{NWPU-RESISC45}

Cheng et al.~\cite{cheng2017nwpu} introduced NWPU-RESISC45 as a remote
sensing scene classification benchmark with 45 classes and 31,500 images,
including 700 images per class. The present experiments use a 20\% training
pool and an 80\% held-out test set. The training pool is divided into a
calibration-training subset and a validation subset before final fitting.
Table~\ref{tab:partitions} gives the retained image counts.

The outer partition seed is 42 and the within-pool validation partition seed
is 4242. These fixed partition seeds are distinct from the seeds used for
training repetitions. The manifest construction identified eight groups of
same-class images with identical pixel content. Each such group is kept
within one subset. All original files are retained, so the manifest contains
31,500 files representing 31,492 unique pixel images. This check addresses
the identified exact duplicates; it does not establish the absence of
near-duplicate images or geographic dependence.

\subsection{Aerial Image Dataset}

The Aerial Image Dataset (AID) was introduced by Xia
et al.~\cite{xia2017aid} for aerial scene classification. It contains 10,000
images from 30 scene classes. The implemented partition uses a 50\% training
pool and a 50\% held-out test set. The pool is divided into 4,000 calibration
images and 1,000 validation images. The outer and within-pool partition seeds
are again 42 and 4242. One same-class exact-duplicate group is confined to
one subset; all 10,000 files are retained, with 9,999 unique pixel images.

\subsection{BigEarthNet-19}

Sumbul et al.~\cite{sumbul2019bigearthnet} introduced BigEarthNet as a
multi-label Sentinel-2 archive. The 19-class nomenclature was described by
Sumbul et al.~\cite{sumbul2021bigearthnetmm}. These dataset references do
not define the number of samples retained by this project's available-data
pipeline. The implemented manifests use the existing partition labels in
BigEarthNet v2.0's \texttt{metadata.parquet}~\cite{clasen2024bigearthnetv2}
and retain patches with the required RGB files. They contain 480,038
samples in total, divided as shown in Table~\ref{tab:partitions}; they are not
a newly generated random 70/15/15 split.

Only Sentinel-2 bands B04, B03 and B02 are used, in red, green and blue order.
When separate band files are read, each band value is divided by 4095,
multiplied by 255, clipped to $[0,255]$ and converted to an unsigned 8-bit
image. The reader also supports previously merged RGB images. The inputs
are therefore three-channel RGB representations, rather than full
multispectral or Sentinel-1/Sentinel-2 inputs. Targets use a fixed ordering
of 19 binary land-cover labels.

\begin{table}[tbp]
 \centering
 \caption{Retained sample counts for the implemented internal protocols.
 ``Pool'' is the union of calibration training and validation for the
 single-label datasets. BigEarthNet keeps training and validation separate.}
 \label{tab:partitions}
 \small
 \setlength{\tabcolsep}{4pt}
 \begin{tabular}{lrrrr}
  \toprule
  Dataset & Training & Validation & Pool & Test\\
  \midrule
  NWPU-RESISC45 & 5,040 & 1,260 & 6,300 & 25,200\\
  AID & 4,000 & 1,000 & 5,000 & 5,000\\
  BigEarthNet-19 & 237,871 & 122,342 & --- & 119,825\\
  \bottomrule
 \end{tabular}
\end{table}

\section{Preprocessing and optimisation}

All models receive $224\times224$ RGB inputs. Evaluation uses deterministic
resizing and ImageNet channel normalisation, with means
$(0.485,0.456,0.406)$ and standard deviations $(0.229,0.224,0.225)$.
The single-label pipeline uses bilinear resizing with antialiasing. During
single-label training, horizontal and vertical flips each have probability
0.5; brightness, contrast and saturation jitter each use magnitude 0.2.
BigEarthNet training uses the resize and two flips, without colour jitter.
Random training augmentations are disabled during validation and testing.

The internal comparison uses minibatches of 32 and the AdamW optimiser,
which was proposed by Loshchilov and Hutter~\cite{loshchilov2019adamw}.
The initial learning rate is $3\times10^{-4}$ and weight decay is
$10^{-4}$. A cosine schedule uses an annealing horizon $T_{\max}=30$ epochs and a
minimum learning rate of zero. The single-label pipeline uses AdamW moment
coefficients $(0.9,0.999)$ and numerical stabiliser $10^{-8}$. It applies
cross-entropy without label smoothing. BigEarthNet uses binary
cross-entropy with logits. The backbone is fine-tuned together with the
classification and graph modules.

The recorded principal runs use 32-bit floating-point computation without
automatic mixed precision (AMP). TensorFloat-32 (TF32) settings are recorded
separately: the single-label pipeline disables TF32 for matrix multiplication
and cuDNN, whereas the documented BigEarthNet soft-seed42 environment disables
matrix-multiplication TF32 and retains cuDNN TF32. These settings describe the
respective pipelines; they do not establish that every historical run used
an identical software or hardware environment.

\section{Model selection and final fitting}

\subsection{Single-label calibration and refitting}

For AID and NWPU-RESISC45, each model first trains for at most 30 calibration
epochs on the internal training subset. Validation overall accuracy selects
the epoch count. The checkpoint is updated only on a strict improvement,
so the earliest occurrence is retained when validation accuracy ties.

After selection, the model is reinitialised from the common pretrained
backbone checkpoint and fitted on the full training pool for the selected
number of epochs. This refit is not continued training from the calibration
checkpoint. The cosine schedule retains its annealing horizon $T_{\max}=30$ epochs;
it is not rescaled to the selected refit duration. The final refit checkpoint
is evaluated on the held-out test set. Thus, a calibration validation score
and a refit test score can come from different trained weights. A reversal
between their rankings alone does not demonstrate overfitting of a single
checkpoint.

\subsection{Multi-label validation-based checkpoint selection}

BigEarthNet uses its training partition for optimisation and its validation
partition for checkpoint selection. The planned training budget is 30 epochs,
and the best checkpoint is updated when validation mean average precision
(mAP) strictly improves. The selected checkpoint is then evaluated on the
fixed test partition. Training and validation are not merged for refitting.
The achieved final epoch and exit status are reported per run, since some
historical runs ended before the planned budget.

The F1 decision threshold is fixed at 0.5. Neither checkpoint choice nor
threshold choice uses the test labels. A higher test F1 score does not
replace mAP as the preselected primary multi-label metric.

\section{Classification metrics}

\subsection{Single-label metrics}

Let $M$ be the number of test images, and let $y_i$ and $\hat{y}_i$ be the
true and predicted class of image $i$. The predicted class is the largest-logit
class. Overall accuracy (OA) is
\begin{equation}
 \operatorname{OA}=\frac{1}{M}\sum_{i=1}^{M}
                       \mathbf{1}[\hat{y}_i=y_i].
 \label{eq:oa}
\end{equation}
For class $c$, let $TP_c$, $FP_c$ and $FN_c$ denote true positives, false
positives and false negatives, respectively. The per-class and macro F1
scores are
\begin{equation}
 F1_c=\frac{2TP_c}{2TP_c+FP_c+FN_c},
 \qquad
 F1_{\mathrm{macro}}=\frac{1}{C}\sum_{c=1}^{C}F1_c.
 \label{eq:macro-f1}
\end{equation}
A zero denominator gives an F1 score of zero in the implemented metric.
Per-class accuracy is $TP_c/(TP_c+FN_c)$ when the class has test support.
The confusion matrix uses true classes as rows and predicted classes as
columns. Scores are displayed as percentages in result tables.

\subsection{Multi-label metrics}

Sigmoid probabilities rank the test images independently for each label.
Let $AP_c$ denote class $c$'s average precision, calculated from that ranking
using the recorded evaluation implementation. For $C_+$ classes with at
least one positive test example, mAP is
\begin{equation}
 \operatorname{mAP}=\frac{1}{C_+}
                 \sum_{c:\,TP_c+FN_c>0}AP_c.
 \label{eq:map}
\end{equation}
The intended full BigEarthNet-19 evaluation covers all 19 labels; absent
classes in small engineering subsets are not grounds for treating subset
mAP as a benchmark result. The class-wise AP vector is retained alongside
the aggregate score.

Binary predictions use $\hat{y}_{ic}=\mathbf{1}[p_{ic}\geq0.5]$, where
$p_{ic}$ is image $i$'s sigmoid probability for class $c$. Macro-F1 uses
Equation~\eqref{eq:macro-f1}. Micro-F1 aggregates counts over all labels:
\begin{equation}
 F1_{\mathrm{micro}}=
 \frac{2\sum_c TP_c}
 {2\sum_c TP_c+\sum_c FP_c+\sum_c FN_c}.
 \label{eq:micro-f1}
\end{equation}
Average precision measures ranking, whereas F1 also depends on the fixed
decision threshold. Their changes can therefore have different directions.

\section{Internal comparisons and repetition statistics}

The principal architecture comparison contains the shared Swin-T baseline,
fixed-partition GAViT, and local-prototype soft-assignment GAViT. Within a
dataset, comparisons retain the same recorded data partitions, label order,
input representation and model-selection rule. The soft variant changes
both region aggregation and the feedback assignment. Accordingly, its
comparison with the fixed variant does not separately isolate the effects
of membership softness, the projection parameters and the feedback map.

Let $v_{ms}$ be the test metric of method $m$ in training repetition $s$, and
let $R$ be the number of eligible repetitions in a stated comparison group.
The reported mean $\bar{v}_m$ and sample standard deviation $s_m$ are
\begin{equation}
 \bar{v}_m=\frac{1}{R}\sum_{s=1}^{R}v_{ms},
 \qquad
 s_m=\sqrt{\frac{1}{R-1}\sum_{s=1}^{R}
                    (v_{ms}-\bar{v}_m)^2},\quad R>1.
 \label{eq:repeat-statistics}
\end{equation}
No standard deviation is assigned to a single run. A paired difference is
computed within the same seed before its mean and standard deviation are
calculated. Differences between percentage scores are expressed in
percentage points (pp).

Each result table states its actual seed set and completion criteria.
Historical runs with different code identities, devices, training durations
or evidence completeness are disclosed separately. They are not relabelled
as fully controlled repetitions. All prespecified runs, including negative
results and interrupted runs, remain in the experiment record. A subset
prepared for a progress report does not replace the complete experiment
record or justify discarding unfavourable runs.

Training repetitions share the frozen test set. They measure variation
conditional on that partition and do not create additional independent test
images or reproduce a repeated-random-split protocol. Means and standard
deviations alone do not establish statistical significance or equivalence.
No significance-test outcome is asserted in this partial draft; the final
test specification must match the completed paired comparisons and their
available sample size.

\section{Efficiency and qualitative analysis}

Efficiency reporting separates parameter count, model-forward latency,
throughput, GPU memory and total experiment duration. A model-forward
benchmark uses the same device, input shape, batch size, numerical precision,
warm-up and synchronisation procedure for all compared models. The archived
AID forward benchmark uses GPU-resident synthetic inputs, 20 warm-up passes
and five groups of 50 timed passes. It excludes image reading, preprocessing
and host-to-device transfer. Its throughput therefore describes the measured
forward path, rather than the complete dataset pipeline. Repeated timing
groups are not independent training seeds, and runs made on different GPU
models are not direct evidence of a model's relative execution speed.

Qualitative analysis distinguishes region assignment, explicit spatial and
feature edges, within-region pooling weights, and graph-attention weights.
These visual quantities represent different operations. Feature similarity
and graph attention are not interchangeable, and neither alone proves a
human-semantic relationship or a causal source of classification improvement.
Anchor locations are distinguished from region centroids. Successes,
introduced errors and shared failures are retained within each declared
selection set. The visualised split and checkpoint identity accompany the
figures, since calibration examples cannot substitute for refit test
performance.

\section{Reproducibility and reporting boundaries}

Run records bind the data manifest, pretrained weight identity, model
configuration, source revision, training seed, numerical settings and
checkpoint to their logs and evaluation outputs. Best validation epoch and
last achieved epoch are recorded separately. A model option in the current
source code is not, by itself, evidence that the corresponding formal
experiment was completed.

Earlier NWPU experiments used a different data split, and earlier graph
variants used different connectivity or integration rules. Those results
are not pooled with the present 20\%-training-pool refit protocol. Likewise,
the historical $3\times3$ BigEarthNet experiment uses a separately frozen
configuration; it does not change the $K=16$ definition of the principal
model in Chapter~\ref{chap:method}. External-method reproduction, final
eligible seed groups, inferential statistics and complete cross-model
efficiency tables must be resolved against their own evidence before the
thesis makes a final comparative claim.

\clearpage
\renewcommand{\bibname}{References}
\bibliographystyle{IEEEtran}
\bibliography{references}
\end{document}
